\section{Precedence dan Associativity}
\label{sec:ch09-precedence}

\subsection{Precedence}

Precedence menentukan operator mana yang diikat lebih kuat. Pada
ekspresi \(\code{id} + \code{id} * \code{id}\), perkalian memiliki
precedence lebih tinggi daripada penjumlahan, sehingga pohonnya
memiliki \code{*} lebih dalam daripada \code{+}.

\begin{table}[H]
\centering
\caption{Tingkat precedence operator Subset-C}
\label{tab:ch09-precedence}
\begin{tabular}{@{}c l l@{}}
\toprule
\textbf{Tingkat} & \textbf{Operator} & \textbf{Associativity} \\
\midrule
1 (terendah) & \code{||} & kiri \\
2 & \code{\&\&} & kiri \\
3 & \code{==}, \code{!=} & kiri \\
4 & \code{<}, \code{<=}, \code{>}, \code{>=} & kiri \\
5 & \code{+}, \code{-} & kiri \\
6 & \code{*}, \code{/}, \code{\%} & kiri \\
7 & \code{!}, unary \code{-}, unary \code{+} & kanan \\
8 (tertinggi) & \code{()}, call, literal, identifier & --- \\
\bottomrule
\end{tabular}
\end{table}

\subsection{Associativity}

Associativity menentukan pengelompokan operator dengan precedence
sama. Operator kiri-asosiatif dikelompokkan dari kiri:
\[
a - b - c = (a - b) - c.
\]
Operator kanan-asosiatif dikelompokkan dari kanan:
\[
a = b = c = a = (b = c).
\]
Operator non-asosiatif tidak boleh dirangkai, misalnya
\(a < b < c\) harus ditulis dengan tanda kurung.

\begin{figure}[H]
\centering
\begin{tikzpicture}[node distance=5mm and 6mm]
  \node[phase, text width=2.4cm] (logor) {logor};
  \node[phase, text width=2.4cm, below=of logor] (logand) {logand};
  \node[phase, text width=2.4cm, below=of logand] (eq) {equality};
  \node[phase, text width=2.4cm, below=of eq] (rel) {relation};
  \node[phase, text width=2.4cm, below=of rel] (add) {additive};
  \node[phase, text width=2.4cm, below=of add] (term) {term};
  \node[phase, text width=2.4cm, below=of term] (unary) {unary};
  \node[phase, text width=2.4cm, below=of unary] (primary) {primary};
  \draw[arrow] (logor) -- (logand);
  \draw[arrow] (logand) -- (eq);
  \draw[arrow] (eq) -- (rel);
  \draw[arrow] (rel) -- (add);
  \draw[arrow] (add) -- (term);
  \draw[arrow] (term) -- (unary);
  \draw[arrow] (unary) -- (primary);
\end{tikzpicture}
\caption{Tingkat nonterminal ekspresi: semakin ke bawah, semakin tinggi precedence}
\label{fig:ch09-levels}
\end{figure}

\subsection{Menyatakan Precedence dalam Grammar}

Precedence dinyatakan dengan membuat satu nonterminal per tingkat.
Aturan tingkat bawah muncul sebagai anak dari tingkat atas. Dengan
cara ini, pohon urai otomatis mencerminkan prioritas tanpa aturan
tambahan.

Sebagai contoh, aturan
\[
\text{additive} \to \text{term}\ \{(\text{"+"} \mid \text{"-"})\
\text{term}\}
\]
menempatkan \code{term} di bawah \code{additive}. Karena
\code{term} menangani perkalian, perkalian selalu lebih dalam
daripada penjumlahan.

